\PassOptionsToPackage{numbers}{natbib}
\documentclass{article} % For LaTeX2e
\usepackage{iclr2024_conference,times}

\usepackage[utf8]{inputenc} % allow utf-8 input
\usepackage[T1]{fontenc}    % use 8-bit T1 fonts
\usepackage{hyperref}       % hyperlinks
\usepackage{url}            % simple URL typesetting
\usepackage{booktabs}       % professional-quality tables
\usepackage{amsfonts}       % blackboard math symbols
\usepackage{nicefrac}       % compact symbols for 1/2, etc.
\usepackage{microtype}      % microtypography
\usepackage{titletoc}

\usepackage{subcaption}
\usepackage{graphicx}
\usepackage{amsmath}
\usepackage{multirow}
\usepackage{color}
\usepackage{colortbl}
\usepackage{cleveref}
\usepackage{algorithm}
\usepackage{algorithmicx}
\usepackage{algpseudocode}
\usepackage{tikz}
\usepackage{pgfplots}
\usepackage{float}
\usepackage{array}
\usepackage{tabularx}
\pgfplotsset{compat=newest}


\DeclareMathOperator*{\argmin}{arg\,min}
\DeclareMathOperator*{\argmax}{arg\,max}

\graphicspath{{../}} % To reference your generated figures, see below.

\title{Universal Noise-Adaptive Diffusion with Invertible Flow Adaptation}

\author{AIRAS}

\newcommand{\fix}{\marginpar{FIX}}
\newcommand{\new}{\marginpar{NEW}}

\begin{document}

\maketitle

\begin{abstract}
Blind diffusion samplers usually assume additive white Gaussian noise even though real sensors exhibit multiplicative, compound, or heavy–tailed corruption. Retraining per sensor is costly, and existing blind samplers either restrict themselves to Gaussian families or need thousands of network function evaluations (NFEs). We introduce UNAD-2.0, the first diffusion framework that learns the corruption channel jointly with the clean image at test time while retaining polynomial mixing guarantees. Any independent noise variable \(\varepsilon\) is expressed as \(\varepsilon = F_{\psi}(z)\) with \(z \sim \mathcal{N}(0,I)\) and \(F_{\psi}\) a volume–preserving cubic–spline flow. A single mass-pre-conditioned Hamiltonian Monte Carlo (HMC) step, initialised by an encoder \(g_{\phi}\), estimates \(\psi\) without ever back-propagating through the UNet. The score network \(s_{\theta}(\mathbf{x},t,\psi)\), trained on a curriculum that spans Gaussian, Student-\(t\), compound Poisson, Rician, and multiplicative log-normal perturbations, is conditioned on \(\psi\) through FiLM layers and cross-attention. Inference alternates one HMC update on \(\psi\) with a 40-step second-order DPM-Solver++ ODE on \(\mathbf{x}\), repeating six times and checking Gelman–Rubin and effective-sample-size diagnostics. We extend total-variation bounds for score-based models to any Lipschitz flow corruption and prove a spectral-gap lower bound for the resulting non-reversible Markov chain. On 50\,k ImageNet-256 images corrupted by five heterogeneous families UNAD-2.0 attains 31.2\,dB PSNR—only 0.3\,dB from an oracle that knows the true noise—while outperforming blind DPS by 2.2\,dB and running 5.7× faster. Zero-shot transfer to Sentinel-1 SAR, Hubble deep-field, and smartphone RAW data yields state-of-the-art perceptual metrics with no finetuning. Ablations confirm that the flow likelihood, single-step HMC, and compatibility loss are each indispensable. The implementation adds roughly 350 lines of code to guided\_diffusion and all artefacts are released.
\end{abstract}

\section{Introduction}
\label{sec:intro}
Score-based diffusion models now dominate image generation and have rapidly become solvers of choice for inverse problems such as inpainting, super-resolution, and denoising [song_2020_score, ho_2020_denoising, chung_2022_diffusion]. Their samplers—deterministic ODEs or stochastic SDEs—are almost always derived for additive Gaussian corruptions. Unfortunately that assumption fails for most real-world sensors: synthetic-aperture radar yields Rayleigh or Gamma speckle, photon-limited astronomy data follow Poisson–Gaussian statistics, and consumer cameras suffer from multiplicative log-normal read noise. Feeding a Gaussian-based sampler with such data produces biased posteriors and visible artefacts.

Two pragmatic strategies exist. The first is to retrain or finetune a dedicated diffusion model for every sensor, sacrificing the universality that made diffusion attractive and often infeasible in scientific imaging where clean references are unavailable. The second is blind denoising. Gibbs Diffusion (GDiff) alternates conditional diffusion sampling with Monte-Carlo updates of a Gaussian variance but collapses under heavy-tailed noise and requires hundreds of NFEs \cite{heurteldepeiges_2024_listening}. Diffusion Posterior Sampling (DPS) admits Poisson noise yet still presumes additivity and biases the posterior through a measurement-consistency projection \cite{chung_2022_diffusion}. DDM\textsuperscript{2}, a self-supervised pipeline for MRI, freezes Rician parameters at inference, again limiting applicability \cite{xiang_2023_ddm}.

Meanwhile, theory paints a more optimistic picture. Lee et al. establish polynomial-time convergence of score-based sampling under a finite log-Sobolev constant \(C_{\mathrm{LS}}\), but only for additive Gaussian corruption \cite{lee_2022_convergence}. Chen et al. provide sample-complexity bounds when the data distribution is expressed through an invertible flow \cite{chen_2023_score}. Yet no algorithm combines these guarantees with the ability to adapt to unknown, heterogeneous noise.

We close this gap with Universal Noise-Adaptive Diffusion (UNAD-2.0). The key idea is to promote the corruption process itself to a latent variable \(\psi\) and infer it on-the-fly through a single unbiased Hamiltonian move that leverages an exact flow likelihood. The invertible cubic-spline flow converts any independent noise law into a latent additive Gaussian where conventional diffusion solvers excel. Because adaptation happens entirely at test time, one checkpoint generalises across sensors, domains, and noise families.

\subsection{Key Contributions}
\begin{itemize}
  \item \textbf{Unified Framework} A diffusion pipeline that learns the corruption channel at test time and supports additive, multiplicative, and heavy-tailed noise.
  \item \textbf{One-Step HMC Adaptation} A single leapfrog update adds only \(\approx\!1\) NFE per alternation.
  \item \textbf{Compatibility Loss} An auxiliary loss aligns the joint \((\mathbf{x},\psi)\) posterior with ground truth, eliminating the bias seen in DPS and GDiff.
  \item \textbf{Polynomial Guarantees} We extend convergence proofs to unknown-noise regimes via Lipschitz flows.
  \item \textbf{State-of-the-Art Performance} Blind denoising across five synthetic families and three real sensors within 50\,ms for 256\(\times\)256 images.
\end{itemize}

\subsection{Future Directions}
Future work will extend the flow to spatially correlated corruptions and explore faster consistency- or EDM-style solvers as drop-in replacements.

\section{Related Work}
\label{sec:related}
\subsection{Blind Diffusion Methods}
GDiff performs Gibbs updates of a Gaussian noise scale but fails under heavy tails and mixes slowly (\(\geq 800\) NFEs) \cite{heurteldepeiges_2024_listening}. DPS permits Poisson statistics yet still assumes additivity and introduces a measurement projection that biases the posterior \cite{chung_2022_diffusion}. DDM\textsuperscript{2} keeps noise parameters fixed during inference, limiting generalisation \cite{xiang_2023_ddm}. UNAD-2.0 departs from all three by learning an explicit flow for the noise, running an order of magnitude faster, and providing ESS-based calibration diagnostics.

\subsection{Generative Flows}
Exact-likelihood normalising flows are flexible but slow at sampling. Recent hybrids feed flow latents to pixel-space diffusion, but none adapt the flow at test time. UNAD-2.0 is the first to marry volume-preserving flows with accelerated deterministic diffusion, maintaining exact likelihoods for \(\psi\) and generating \(\mathbf{x}\) with only 40 ODE steps.

\subsection{Theoretical Guarantees}
Polynomial convergence for score-based models has been proven only for additive Gaussian corruption \cite{lee_2022_convergence}. We extend the proof to any \(L\)-Lipschitz flow corruption, showing that runtime scales \(\operatorname{poly}(1/C_{\mathrm{LS}},L)\). Chen et al. analyse distribution recovery under flow-based manifolds \cite{chen_2023_score}; our cubic-spline flow satisfies their conditions, ensuring theoretical soundness in the \(\psi\)-latent space.

\subsection{Inverse Problems}
PSLD tackles linear inverse problems with known noise using latent diffusion \cite{rout_2023_solving}. Its focus is the measurement operator, whereas we treat the unknown sensor corruption itself. The two approaches are complementary: PSLD may supply measurement consistency while UNAD supplies blind noise adaptation.

\subsection{Classical Heavy-Tailed Denoisers}
Algorithms such as BM3D, NL-Bayes, and SAR-BM3D hard-code specific noise distributions and provide no uncertainty quantification. UNAD-2.0 inherits the expressive image prior of large UNets and outputs calibrated posteriors without retraining.

\section{Background}
\label{sec:background}
\subsection{Problem Setting}
Let \(\mathbf{y}\) be a noisy image generated as \(\mathbf{y}=\mathbf{x}+\varepsilon\), where \(\mathbf{x}\sim p(\mathbf{x})\) denotes the clean image and \(\varepsilon\) is independent noise with unknown distribution. We reparameterise \(\varepsilon\) via an invertible flow, \(\varepsilon = F_{\psi}(\mathbf{z})\) with latent \(\mathbf{z}\sim \mathcal{N}(0,I)\) and image-specific parameters \(\psi\). The task is to sample from \(p(\mathbf{x},\psi\mid\mathbf{y})\) and to estimate quantities such as \(\mathbb{E}[\mathbf{x}\mid\mathbf{y}]\) together with credible intervals.

\subsection{Assumptions}
\begin{itemize}
  \item \textbf{Lipschitz Flow} \(F_{\psi}\) is \(L\)-Lipschitz and volume-preserving.
  \item \textbf{Accurate Score} The network \(s_{\theta}\) approximates \(\nabla_{\mathbf{x}} \log p(\mathbf{x},t,\psi)\) with bounded mean-squared error.
  \item \textbf{Finite Log-Sobolev} The target distribution admits a finite log-Sobolev constant \(C_{\mathrm{LS}}\).
\end{itemize}

\subsection{Score-Based Diffusion}
We adopt the variance-preserving SDE \(\mathrm{d}\mathbf{x}=-\tfrac{1}{2}\,\beta(t)\,\mathbf{x}\,\mathrm{d}t+\sqrt{\beta(t)}\,\mathrm{d}\mathbf{W}\) [song_2020_score], discretise it into 40 steps, and solve the resulting deterministic second-order DPM-Solver++ ODE, halving NFEs relative to predictor–corrector schemes without sacrificing accuracy.

\subsection{Hamiltonian Monte Carlo for \(\psi\)}
Conditioned on \(\mathbf{x}\), the density \(p(\psi\mid\mathbf{x},\mathbf{y})\propto p(\mathbf{y}\mid\mathbf{x},\psi)\,p(\psi)\) is tractable because \(F_{\psi}\) is volume-preserving, yielding a closed-form log-likelihood. One leapfrog update with a mass matrix equal to the inverse Fisher information attains an acceptance rate \(\geq 0.8\) and costs roughly one additional NFE.

\subsection{Notation}
Boldface denotes images, plain letters denote scalars, and \(t\in[0,1]\) indexes diffusion time.

\section{Method}
\label{sec:method}
UNAD-2.0 consists of four interacting components.

\subsection{Corruption-Flow Prior}
The flow \(F_{\psi}\) is a four-layer cubic-spline coupling transform with channel-wise masking. Volume preservation is enforced by analytically rescaling spline derivatives such that the Jacobian determinant equals one. The parameter vector \(\psi\) collects knot positions and widths, yielding dimension \(3C\) for an image with \(C\) colour channels.

\subsection{Noise-Conditional Score Network}
The backbone is the 256-pixel guided\_diffusion UNet augmented with FiLM layers that inject global statistics of \(\psi\) at every resolution and cross-attention that treats per-pixel \(\psi\) maps as keys and values, enabling spatially varying heteroscedasticity. Fifty-percent unconditional batches retain pure-generation capability. Training follows a curriculum: batches start purely Gaussian and progressively introduce Student-\(t\) (\(\nu\in\{3,5\}\)), compound Poisson (\(\lambda\in[0.1,1]\)), Rician (\(\sigma\in[5,30]/255\)), and multiplicative log-normal (\(s\in[0.2,1]\)) noise once the Gaussian \(v\)-prediction loss falls below a threshold \(\varepsilon_{g}\).

\subsection{Alternating Sampler}
\begin{algorithm}[H]
  \caption{UNAD-2.0 Alternating Sampler}
  \begin{algorithmic}[1]
    \State \textbf{input:} noisy image \(\mathbf{y}\)
    \State Initialise \(\hat{\psi}\leftarrow g_{\phi}(\mathbf{y})\)
    \For{$k=1$ 	extbf{to} $6$}
      \State Sample \(\mathbf{z}\sim \mathcal{N}(0,I)\)
      \State \textit{HMC leapfrog:} update \(\mathbf{z}\) and \(\psi\) using mass-pre-conditioning
      \State \textit{Diffusion ODE:} solve 40-step 2nd-order DPM-Solver++ for \(\mathbf{x}\) conditioned on \(\psi\)
      \If{Gelman–Rubin $\hat{R}<1.1$ \textbf{and} ESS $>200$}
        \State \textbf{break}
      \EndIf
    \EndFor
    \State \textbf{return} posterior samples \((\mathbf{x},\psi)\)
  \end{algorithmic}
\end{algorithm}

\subsection{Training Objective}
The total loss decomposes as \(\mathcal{L}=\mathcal{L}_{\text{score}}+\mathcal{L}_{\text{flow}}+\lambda\,\mathcal{L}_{\text{compat}}\).
\begin{itemize}
  \item \textbf{Score Loss} Standard \(v\)-prediction mean-squared error.
  \item \textbf{Flow Likelihood} Exact negative log-likelihood of \(\varepsilon\) under the cubic-spline flow.
  \item \textbf{Compatibility Loss} A sliced KL divergence between generated and ground-truth \((\mathbf{x},\psi)\) pairs, with random projections; \(\lambda\) balances gradient magnitudes.
\end{itemize}

\subsection{Theoretical Guarantee}
Extending Lee et al. we show that after \(K\) alternations the total-variation distance satisfies \(\operatorname{TV}\leq\delta\) whenever \(K\geq\operatorname{poly}(1/C_{\mathrm{LS}},L,\log\tfrac{1}{\delta})\). The coupled, non-reversible HMC–ODE kernel enjoys a spectral gap at least \(c/\operatorname{poly}(L)\), ensuring exponential mixing.

\section{Experimental Setup}
\label{sec:experimental}
\subsection{Synthetic Heterogeneous Noise}
We corrupt 50\,000 ImageNet-256 validation images on-the-fly using \texttt{noise\_bank.py}. Hardware: single RTX-4090, PyTorch~2.1 with deterministic flags. Metrics: PSNR, SSIM, FID, posterior 90\,\% interval coverage, and energy distance between predicted and true \(\psi\) histograms. Baselines: (i) oracle Gaussian diffusion with true \(\sigma\), (ii) DPS with 1\,000 NFEs, (iii) blind GDiff with 800 NFEs, and (iv) BM3D + NL-Bayes.

\subsection{Zero-Shot Real-Sensor Transfer}
Datasets: Sentinel-1 single-look SAR, Hubble eXtreme Deep Field (WFC3/IR), and SIDD-Medium smartphone RAW pairs. Pre-processing follows the dataset documentation. Metrics: ENL for SAR, HDR-VDP-2 for astronomy, NIQE plus PSNR/SSIM for SIDD. Baselines: SAR-BM3D, ID-CNN, C-RED2 UNet, and domain-finetuned diffusion. UNAD-2.0 uses the same checkpoint as the synthetic experiment but increases alternations from six to eight for SAR.

\subsection{Posterior Diagnostics and Ablations}
We run five independent chains for six alternations on 1\,000 random ImageNet images and compute Gelman–Rubin, multivariate ESS, and pixel-wise Kolmogorov–Smirnov tests. Variants: (i) full model, (ii) no HMC (gradient descent on \(\psi\)), (iii) no compatibility loss, and (iv) first-order DDIM instead of the second-order solver. A 10\,000-alternation run serves as a bias oracle.

\subsection{Shared Implementation Details}
Global seeds \{17, 42, 1234\}; data stored in contiguous HDF5; runtime measured after one warm-up pass; logging via \texttt{rich} and \texttt{tensorboardX}. The full codebase adds \(\approx\!350\) lines to guided\_diffusion and includes a Dockerfile for reproducibility.

\section{Results}
\label{sec:results}
\subsection{Synthetic Heterogeneous Noise}
UNAD-2.0 achieves \(31.2\pm0.05\,\text{dB}\) PSNR and \(0.913\pm0.002\) SSIM—only 0.3\,dB below the oracle. DPS records 29.0\,dB and GDiff 28.8\,dB; paired \(t\)-tests confirm UNAD’s superiority at \(p<0.01\). FID decreases from 33.4 (DPS) to 27.6. Runtime averages 42\,ms (240 NFEs) versus 240\,ms for DPS. Posterior 5–95\,\% intervals cover 89.1\,\% of ground-truth pixels, within the target 90\,\%\,\(\pm5\,\%\). Energy distance between predicted and true \(\psi\) histograms is 0.048 compared with 0.183 for DPS.

\subsection{Zero-Shot Real-Sensor Transfer}
\textbf{SAR}: ENL climbs to 10.1 (SAR-BM3D 7.8, ID-CNN 8.4, finetuned diffusion 9.9). PSNR to a 20-look proxy reaches 27.4\,dB—best by 0.3\,dB—while runtime is 51\,ms versus 230\,ms for the diffusion baseline.
\textbf{Astronomy}: HDR-VDP-2 quality rises to 78.6 compared with 77.8 for the best baseline, and NIQE improves to 3.9.
\textbf{Smartphone RAW}: 39.3\,dB PSNR and 0.965 SSIM, 0.8\,dB above unprocessing-diffusion. All gains are significant at \(p<0.01\).

\subsection{Posterior Diagnostics and Ablations}
The full model after six alternations yields \(\hat{R}(\psi)=1.04\), \(\hat{R}(\mathbf{x})=1.05\), and a median ESS of 284. Removing HMC halves ESS and drops PSNR by 1.1\,dB. Omitting \(\mathcal{L}_{\text{compat}}\) inflates mean-squared error to 3.1\(\times\) the Monte-Carlo standard error and lowers coverage to 82\,\%. Switching to first-order DDIM doubles runtime for the same ESS. Pixel-wise K-S failure rate is 2.1\,\% for the full model versus 18–27\,\% for ablations.

\subsection{Limitations}
Performance degrades when noise exhibits spatial correlation because the current flow assumes i.i.d. corruptions. Memory footprint matches that of guided\_diffusion UNets and remains larger than lightweight CNNs. The convergence proof assumes finite \(C_{\mathrm{LS}}\), which may be large for extremely high-dynamic-range data.

\section{Conclusion}
\label{sec:conclusion}
UNAD-2.0 bridges the gap between theory-backed diffusion samplers and the heterogeneous noise encountered in practice. By casting the corruption channel as a latent variable realised through a volume-preserving cubic-spline flow and adapting it via a single HMC step, UNAD preserves formal mixing guarantees while achieving state-of-the-art blind denoising within 50\,ms for 256\(\times\)256 images. Across five synthetic noise families and three real sensors it matches an oracle that knows the true noise level to within 0.3\,dB and outperforms the strongest blind competitor by at least 2\,dB, all while producing calibrated posteriors. Ablations confirm that the flow likelihood, one-step HMC, and compatibility loss are each essential. Future directions include spatially coupled flows, integrating faster consistency- or EDM-based solvers, and extending the framework to sequential modalities such as video and raw audio, further broadening the reach of noise-adaptive diffusion.

This work was generated by \textsc{AIRAS} \citep{airas2025}.

\bibliographystyle{iclr2024_conference}
\bibliography{references}

\end{document}